\documentclass[12pt,a4paper]{article}
\usepackage[utf8]{inputenc}
\usepackage[T1]{fontenc}
\usepackage[italian]{babel}
\usepackage{libertinus}
\usepackage{amssymb}
\usepackage[a4paper,top=25mm,bottom=25mm,left=22mm,right=22mm,headheight=15pt]{geometry}
\usepackage{graphicx}
\graphicspath{{../../../assets/}{../../assets/}{../assets/}{assets/}}
\usepackage[table]{xcolor}
\usepackage{tabularx,array}
\usepackage{titlesec}
\usepackage{titletoc}
\usepackage{microtype}
\usepackage{fancyhdr}
\usepackage{enumitem}
\usepackage[hidelinks]{hyperref}

% ---------- Dati del verbale (da impostare una sola volta) ----------
\newcommand{\vnum}{\ph{n}}                     % numero progressivo del verbale interno
\newcommand{\vdata}{\ph{AAAA-MM-GG}}           % data della riunione, formato ISO
\newcommand{\vdataestesa}{\ph{GG mese AAAA}}   % data della riunione per esteso (es. 19 novembre 2026)

% Palette colori del gruppo Bending Forks
\definecolor{bfblue}{HTML}{1B3F9E}
\definecolor{bfrow}{HTML}{F2F5FC}
\definecolor{bfline}{HTML}{C9D1E6}
\definecolor{bfhead}{HTML}{2E6BE6}

% Macro di formattazione interna Bending Forks
\newcommand{\ph}[1]{\textcolor{red}{[#1]}}                        % segnaposto da completare
\newcommand{\gl}[1]{\emph{#1}\textsubscript{\scriptsize G}}       % termine del Glossario
\newcommand{\thc}[1]{\textcolor{white}{\textbf{#1}}}              % cella di intestazione delle tabelle
\newcommand{\rl}{\arrayrulecolor{bfline}\hline}                   % riga di separazione delle tabelle
\newcommand{\presente}{$\blacksquare$ & $\square$}                % colonne Presente / Assente
\newcommand{\assente}{$\square$ & $\blacksquare$}
\newcommand{\tightlist}{\setlength{\itemsep}{1pt}\setlength{\parskip}{0pt}}

\setlist{topsep=2pt}
\setlist[itemize]{leftmargin=1.3em,label=\textbullet}
\setlist[enumerate]{leftmargin=1.6em}
\setlength{\parindent}{0pt}
\sloppy
\setlength{\emergencystretch}{2em}
\setlength{\parskip}{5pt}
\renewcommand{\arraystretch}{1.45}
\setcounter{secnumdepth}{3}
\setcounter{tocdepth}{3}

% Formattazione titoli sezioni
\titleformat{\section}{\LARGE\bfseries\color{bfblue}\raggedright}{\thesection.}{0.3em}{}[{\color{bfblue}\vspace{-2pt}\titlerule[0.9pt]}]
\titleformat{\subsection}{\Large\bfseries\color{bfblue}\raggedright}{\thesubsection.}{0.3em}{}
\titleformat{\subsubsection}{\normalsize\bfseries\raggedright}{\thesubsubsection.}{0.3em}{}
\titlespacing*{\section}{0pt}{14pt}{6pt}
\titlespacing*{\subsection}{0pt}{10pt}{4pt}
\titlespacing*{\subsubsection}{0pt}{10pt}{3pt}

% Indice cliccabile con trattini fino al numero di pagina
\contentsmargin{2.2em}
\newcommand{\bfleader}{\titlerule*[0.32em]{\textendash}}
\titlecontents{section}[0pt]{\addvspace{5pt}\large\bfseries}{\thecontentslabel.\ }{}{\hspace{0.4em}\bfleader\contentspage}
\titlecontents{subsection}[1.2em]{\addvspace{0pt}}{\thecontentslabel.\ }{}{\hspace{0.4em}\bfleader\contentspage}
\titlecontents{subsubsection}[2.4em]{\addvspace{0pt}}{\thecontentslabel.\ }{}{\hspace{0.4em}\bfleader\contentspage}

% Intestazione e piè di pagina
\pagestyle{fancy}
\fancyhf{}
\fancyhead[L]{\footnotesize\sffamily\bfseries\color{bfhead}\textls[120]{BENDING FORKS}}
\fancyhead[R]{\footnotesize\sffamily Verbale interno n. \vnum}
\fancyfoot[C]{\small Pagina \thepage}
\renewcommand{\headrulewidth}{0.9pt}

% Colorazione alternata delle righe delle tabelle
\newcommand{\bftable}{\rowcolors{2}{bfrow}{white}}
\usepackage{etoolbox}
\AtBeginEnvironment{tabularx}{\bftable}

\begin{document}

% ---------- Copertina ----------
\begin{titlepage}
\centering
\vspace*{4mm}
\includegraphics[width=105mm]{logo-bending-forks-trasparente.png}\\[4mm]
\underline{bending.forks.swe@gmail.com}

\vspace{42mm}
{\fontsize{24}{28}\selectfont\bfseries Verbale interno n. \vnum\par}
\vspace{4mm}
{\Large\color{bfblue} Riunione del \vdataestesa\par}

\vfill
\begin{tabular}[t]{r|l}
\textbf{Responsabile} & \ph{Nome Cognome} \\
\textbf{Segretario} & \ph{Nome Cognome} \\
\textbf{Verificatore} & \ph{Nome Cognome} \\
\end{tabular}
\hspace{14mm}
\begin{tabular}[t]{r|l}
\textbf{Versione} & 0.1.0 \\
\textbf{Tipo} & Verbale interno \\
\textbf{Destinatari} & \emph{Gruppo Bending Forks} \\
 & \emph{Prof. Tullio Vardanega} \\
 & \emph{Prof. Riccardo Cardin} \\
\end{tabular}

\vspace{12mm}
{\large\emph{\textbf{Abstract dei contenuti:}}}\\[1mm]
\begin{minipage}{150mm}\centering Sintesi e decisioni emerse durante la riunione interna del gruppo del \vdataestesa, riguardante \ph{argomenti principali della riunione}.\end{minipage}
\vspace*{6mm}
\end{titlepage}

% ---------- Indice ----------
\vspace*{2mm}
{\fontsize{22}{26}\selectfont\bfseries Indice\par}
\vspace{6mm}
\makeatletter\@starttoc{toc}\makeatother
\newpage

% ---------- Registro delle versioni ----------
% Una riga per versione, dalla più recente alla più vecchia.
% Versione nel formato X.Y.Z (Norme di progetto, Versionamento dei documenti): la prima stesura
% è 0.1.0, le correzioni aumentano Z (es. 0.1.1), l'approvazione del Responsabile porta a 1.0.0.
% La versione in copertina coincide con quella dell'ultima riga del registro.
% Il Verificatore inserisce il proprio nome nella riga della versione che ha verificato.
\hypertarget{registro-delle-versioni}{%
\section{Registro delle versioni}}

\begin{center}
\begin{tabularx}{\linewidth}{
  >{\hsize=0.55\hsize\raggedright\arraybackslash}X
  >{\hsize=0.75\hsize\raggedright\arraybackslash}X
  >{\hsize=0.85\hsize\raggedright\arraybackslash}X
  >{\hsize=0.85\hsize\raggedright\arraybackslash}X
  >{\hsize=2.00\hsize\raggedright\arraybackslash}X
}
\rowcolor{bfblue}\thc{Versione} & \thc{Data} & \thc{Autore} & \thc{Verificatore} & \thc{Descrizione} \\
0.1.0 & \ph{AAAA-MM-GG} & \ph{Segretario} & \ph{Verificatore} & Redazione del verbale \\ \rl
\end{tabularx}
\end{center}

% ---------- Informazioni generali ----------
\hypertarget{informazioni-generali}{%
\section{Informazioni generali}}

\hypertarget{durata-e-luogo}{%
\subsection{Durata e luogo}}

\begin{itemize}
\tightlist
\item \textbf{Data:} \vdata;
\item \textbf{Inizio:} \ph{HH:MM};
\item \textbf{Fine:} \ph{HH:MM};
\item \textbf{Luogo / modalità:} in remoto, tramite Discord.
\end{itemize}

\hypertarget{partecipanti}{%
\subsection{Partecipanti}}

% Per ogni componente: ruolo in carica nel periodo e \presente oppure \assente.
\begin{center}
\begin{tabularx}{\linewidth}{>{\raggedright\arraybackslash}X>{\centering\arraybackslash}X>{\centering\arraybackslash}p{2.2cm}>{\centering\arraybackslash}p{2.2cm}}
\rowcolor{bfblue}\thc{Nome e cognome} & \thc{Ruolo} & \thc{Presente} & \thc{Assente} \\
Marco Beltrame & \ph{Ruolo} & \presente \\ \rl
Lorenzo Milanese & \ph{Ruolo} & \presente \\ \rl
Filippo Rampazzo & \ph{Ruolo} & \presente \\ \rl
Nicholas Salmaso & \ph{Ruolo} & \presente \\ \rl
Luca Stangherlin & \ph{Ruolo} & \presente \\ \rl
Aldo Toffoli & \ph{Ruolo} & \presente \\ \rl
Matteo Wu & \ph{Ruolo} & \presente \\ \rl
\end{tabularx}
\end{center}

% ---------- Ordine del giorno ----------
\hypertarget{ordine-del-giorno}{%
\section{Ordine del giorno}}

\begin{enumerate}
\tightlist
\item \hyperlink{odg-1}{\ph{Primo punto dell'ordine del giorno}};
\item \hyperlink{odg-2}{\ph{Secondo punto dell'ordine del giorno}};
\item \hyperlink{varie-ed-eventuali}{Varie ed eventuali}.
\end{enumerate}
\newpage

% ---------- Discussione e decisioni ----------
% Una sottosezione per ogni punto dell'ordine del giorno, con lo stesso titolo usato nell'elenco.
% Ogni decisione ha il codice [VI-<n>-<nn>]: <n> è il numero del verbale (\vnum),
% <nn> il numero progressivo della decisione nel verbale, a due cifre (01, 02, ...).
% Motivazione, incaricati e scadenza si scrivono solo se il punto contiene decisioni;
% un punto senza decisioni ha solo la discussione.
\hypertarget{discussione-e-decisioni}{%
\section{Discussione e decisioni}}

\hypertarget{odg-1}{%
\subsection{\ph{Primo punto dell'ordine del giorno}}}

\textbf{Discussione:}\\
\ph{Sintesi della discussione: chi ha esposto cosa, proposte avanzate, dubbi emersi ed esito.}

\textbf{Decisioni prese:}
\begin{itemize}
\tightlist
\item[$\checkmark$] \textbf{[VI-\vnum-01]} \ph{Prima decisione};
\item[$\checkmark$] \textbf{[VI-\vnum-02]} \ph{Seconda decisione}.
\end{itemize}

\textbf{Motivazione:}\\
\ph{Perché il gruppo ha preso queste decisioni.}

\textbf{Incaricati:} \ph{Cognome, Cognome}\\
\textbf{Scadenza:} \ph{AAAA-MM-GG}

\hypertarget{odg-2}{%
\subsection{\ph{Secondo punto dell'ordine del giorno}}}

\textbf{Discussione:}\\
\ph{Sintesi della discussione.}

\textbf{Decisioni prese:}
\begin{itemize}
\tightlist
\item[$\checkmark$] \textbf{[VI-\vnum-03]} \ph{Decisione}.
\end{itemize}

\textbf{Motivazione:}\\
\ph{Perché il gruppo ha preso questa decisione.}

\textbf{Incaricati:} \ph{Cognome}\\
\textbf{Scadenza:} \ph{AAAA-MM-GG}

% Esempio facoltativo: tabella della rotazione dei ruoli, da inserire nel punto in cui si assegnano i ruoli
% e da richiamare nella decisione con \hyperlink{tab-ruoli}{tabella dei ruoli}.
%
% \hypertarget{tab-ruoli}{%
% \begin{center}
% \begin{tabularx}{\linewidth}{>{\raggedright\arraybackslash}X>{\raggedright\arraybackslash}X}
% \rowcolor{bfblue}\thc{Nome e cognome} & \thc{Ruolo} \\
% Marco Beltrame & \ph{Ruolo} \\ \rl
% Lorenzo Milanese & \ph{Ruolo} \\ \rl
% Filippo Rampazzo & \ph{Ruolo} \\ \rl
% Nicholas Salmaso & \ph{Ruolo} \\ \rl
% Luca Stangherlin & \ph{Ruolo} \\ \rl
% Aldo Toffoli & \ph{Ruolo} \\ \rl
% Matteo Wu & \ph{Ruolo} \\ \rl
% \end{tabularx}\\[2pt]
% \textbf{Suddivisione dei ruoli dal \ph{AAAA-MM-GG} al \ph{AAAA-MM-GG}}
% \end{center}}

\hypertarget{varie-ed-eventuali}{%
\subsection{Varie ed eventuali}}

\textbf{Discussione:}\\
\ph{Argomenti emersi fuori dall'ordine del giorno; se non ce ne sono, scrivere ``Nessun argomento ulteriore.''}

\newpage

% ---------- Azioni da intraprendere ----------
% Una riga per ogni attività assegnata durante la riunione. Priorità: Alta, Media o Bassa.
\hypertarget{azioni-da-intraprendere}{%
\section{Azioni da intraprendere}}

\begin{center}
\begin{tabularx}{\linewidth}{>{\hsize=2.0\hsize\raggedright\arraybackslash}X >{\hsize=1.0\hsize\raggedright\arraybackslash}X >{\hsize=0.5\hsize\centering\arraybackslash}X >{\hsize=0.5\hsize\centering\arraybackslash}X}
\rowcolor{bfblue}\thc{Descrizione} & \thc{Incaricati} & \thc{Scadenza} & \thc{Priorità} \\
\ph{Attività da svolgere} & \ph{Cognome} & \ph{AAAA-MM-GG} & \ph{Alta} \\ \rl
\ph{Attività da svolgere} & \ph{Cognome} & \ph{AAAA-MM-GG} & \ph{Media} \\ \rl
\end{tabularx}
\end{center}

% ---------- Prossima riunione ----------
\hypertarget{prossima-riunione}{%
\section{Pianificazione della prossima riunione}}

\begin{itemize}
\tightlist
\item \textbf{Data proposta:} \ph{AAAA-MM-GG};
\item \textbf{Ora prevista:} \ph{HH:MM};
\item \textbf{Luogo / modalità:} in remoto, tramite Discord;
\item \textbf{Ordine del giorno preliminare:} \ph{argomenti previsti}.
\end{itemize}

% ---------- Firme ----------
% Le firme si aggiungono dopo la verifica: per ciascun firmatario, subito prima di \rule,
% inserire la sua immagine della firma presa da assets/ (file firma_<nome>), racchiusa in
% un makebox largo 0pt, come nei verbali precedenti. Finché le immagini mancano,
% il sito considera il verbale non firmato.
\vspace{35pt}
\noindent
\begin{minipage}[t]{0.45\textwidth}
\textbf{Firma Segretario:}\\[18pt]
\rule{55mm}{0.5pt}\\[2pt]
\small \ph{Nome Cognome}
\end{minipage}
\hfill
\begin{minipage}[t]{0.45\textwidth}
\begin{flushright}
\textbf{Firma Verificatore:}\\[18pt]
\rule{55mm}{0.5pt}\\[2pt]
\small \ph{Nome Cognome}
\end{flushright}
\end{minipage}

\end{document}
